\section{Deployment and Operations}
\label{sec:deployment-ops}

A research engine is only as reproducible as the procedure that ships it and the
discipline that keeps it running. This section records how the kernel reaches Cloud
Run, the isolation and credit posture it runs under, and the two standing
operations that keep the system live without supervision: the nightly self-check
loop and the systemic-risk feed. Throughout, the governing constraint is the same
as in the architecture -- we expose nothing whose cost or blast radius we cannot
bound in advance.

\subsection{Deploying the kernel to Cloud Run}
\label{sec:deploy}

The kernel is deployed by \file{cloudrun/deploy.sh}, which builds a temporary
context and runs \code{gcloud run deploy -{}-source}. The build step exists because
the one bundled dataset, \file{G20.xlsx}, lives outside the engine directory and
Docker cannot \code{COPY} from outside its context; the script stages the server,
the R analyses, the web assets, and the three published method packages
(\code{sochcontagion}, \code{contagionchannels}, \code{namh}) into a scratch
directory alongside the dataset, then deploys that. The image is defined by
\file{cloudrun/Dockerfile} from \code{rocker/r-ver:4.4.1}.

The R stack drives most of the image. The \code{igraph} binary links against GLPK,
libxml2 and GMP, so the Dockerfile installs the system libraries
\code{libglpk40}, \code{libxml2} and \code{libgmp10} \emph{before} the R package
install, both so the post-install load check succeeds and because the method needs
them at runtime; the remaining R packages are pure-R or Rcpp and need no system
dependencies. The three published packages are installed from their bundled
tarballs, so a method such as \method{channel\_attribution} or \method{namh\_te}
runs the paper's own code, not a reimplementation, and the build fails loudly if any
package or its bundled data is missing. Node 20 is added last; the server itself is
zero-dependency.

Two deployment facts carry operational weight. First, the model API key is injected,
never baked in. The script resolves \code{GOOGLE\_API\_KEY} from the environment or
from a \code{.env.local} one level above the repository and passes it via
\code{-{}-set-env-vars}; it is never printed and never committed. If the key is
absent the script prints an explicit warning and ships a revision on which
\route{/api/chat} is simply disabled -- a degraded but honest service, not a broken
one. A revision accidentally shipped without the key can be repaired in place with
\code{gcloud run services update -{}-update-env-vars}, no rebuild required.

Second, the deploy flags encode the cost posture. The service is deployed with
\code{-{}-timeout 300}, \code{-{}-min-instances 0}, \code{-{}-max-instances 2} and
\code{-{}-concurrency 16}, alongside the runtime caps \code{COMPUTE\_TIMEOUT\_S=90},
\code{MAX\_CONCURRENT=8} and \code{MAX\_LLM\_PER\_DAY=400}. The \code{-{}-min-instances
0} is the scale-to-zero economics of Section~\ref{sec:kernel}: between requests the
service costs nothing, which is appropriate for occasional, bursty research traffic
and is what makes a publicly reachable compute endpoint affordable to operate
indefinitely. The remaining caps are the credit and abuse controls of
Section~\ref{sec:cost-bound}.

\subsection{Isolation and the Tier-0 hardening posture}
\label{sec:hardening}

On Cloud Run, gVisor is the isolation boundary around the container, so the server's
\code{bwrap} path is skipped by auto-detection and analyses run under the
\code{timeout} fallback (Section~\ref{sec:sandbox}); the container holds no secrets
beyond the injected key and the G20 panel is the only bundled dataset. We are
explicit that this is \code{timeout} mode, not a network-isolated \code{bwrap}
sandbox -- the live \route{/health} says so -- and the no-arbitrary-code guarantee
therefore rests on the parameterised-only registry (Invariant~\ref{inv:no-rce}),
which holds in either mode.

On top of that boundary sits a layer of hardening whose purpose is to make a public
launch safe against credit exhaustion and abuse rather than against code injection.
The kernel applies, before any model call or spawn: a global per-instance
concurrency ceiling (\code{MAX\_CONCURRENT}, shedding excess with \code{503}); a
per-IP per-minute limiter and a per-IP daily quota on the paid routes; a 64\,KB
request-body cap and a message-length cap that reject oversized or over-long inputs
with \code{413} before any Gemini call is made; a slowloris guard on request and
header timeouts; and a stale-cache fallback that serves the last good result,
clearly flagged, when a live run fails rather than returning an error. Adversarial
load testing of this posture -- per-IP bursts, multi-IP fan-out, oversized bodies,
over-length inputs, method-injection probes, and budget exhaustion -- produced no
uncoded error and no \code{500}; every rejection was a coded JSON response. The
binding financial control among these is the daily model budget, treated next.

\subsection{The cost bound on the AI layer}
\label{sec:cost-bound}

The AI routes are the only paid component, and their spend is bounded by a hard,
verifiable ceiling rather than by a pricing estimate. Each instance enforces a
global per-day cap of \code{MAX\_LLM\_PER\_DAY = 400} requests that reach the paid
stage, reset at UTC midnight, on top of the per-IP limits. Because Cloud Run is
deployed with \code{-{}-max-instances 2}, the fleet-wide ceiling is the product:
\[
  \text{paid turns/day} \;\le\; \underbrace{400}_{\text{per instance}}
  \;\times\; \underbrace{2}_{\text{max instances}} \;=\; 800 .
\]
At most 800 paid turns per day can occur regardless of how many distinct IPs attack
the service. A ``turn'' is one chat or research request that reaches the paid stage;
over the cap, the route returns \code{503 daily\_capacity}, never a fabricated
answer. This bound is what makes the launch safe: it is independent of attacker IP
count, it is enforced before the Gemini call so an over-cap request spends nothing,
and it is tunable through the environment. The application-level caps are the active
backstop; a billing-account budget alert is an additional financial tripwire to be
set by the project owner.

\subsection{The governed cloud-capability layer}
\label{sec:capability-layer}

The cost bound of the previous subsection generalises into a small framework, because
the language-model analyst is no longer the only managed-cloud surface the engine can
reach. A governed capability layer in \file{server/upgrade-capabilities.mjs} exposes a
fixed menu of paid Google Cloud services through one route, \route{POST /api/upgrade/run},
whose body names a capability and a set of typed parameters. The route reuses the registry
discipline of Section~\ref{sec:registry-guard}: it matches the name against a fixed
four-entry table, refuses anything else with a coded error, and is fetch-only, so it adds
no new execution surface and the parameterised-only guard is unchanged.

Each capability is disabled unless its own flag, \code{CE\_CAP\_<NAME>}, is set
explicitly, so the resting posture is least-privilege and off. Every paid call then passes
a fixed gate in order, the flag, the typed-parameter check, a prerequisite check where one
applies, and last a spend ceiling, before any request leaves the engine. A rejection at
any stage returns a coded capacity or argument error and spends nothing, never a
fabricated answer, exactly as the analyst cost bound does. The ceiling is one reusable
governor, \method{capBudget}, in the same guards module as the analyst budget; it keeps a
per-UTC-day counter for each named capability, billed in calls or, for embeddings, in
tokens, and it reserves the units \emph{before} the call so that an over-cap request is
refused rather than charged. It generalises the single \code{MAX\_LLM\_PER\_DAY} counter
of Section~\ref{sec:cost-bound} to a family of named counters, and like the core budget of
Section~\ref{sec:core-governance} it is the governor a saturable resource must ship with
before it is reachable. We confirmed the ceiling on the live service: with the
Gemini-on-Vertex flash cap set to one for the test, the second request of the day returned
the coded capacity error rather than a result, and the cap was then restored.

\begin{table}[h]
\centering\small
\begin{tabular}{@{}llll@{}}
\toprule
Capability & Vertex surface & Daily ceiling & Status \\
\midrule
Embeddings        & \code{text-embedding-004}        & $2{,}000{,}000$ tokens          & live \\
Grounded search   & Vertex AI Search                 & $200$ queries                   & live \\
Gemini-on-Vertex  & \code{aiplatform} generate       & $300$ flash / $100$ pro         & live \\
Batch prediction  & \code{batchPredictionJobs}       & $1$ batch, $\le 50{,}000$ items & live \\
\bottomrule
\end{tabular}
\caption{The four governed cloud capabilities, each behind a default-off flag with a
per-UTC-day spend ceiling that trips before the paid call. All four are live end-to-end on
the deployed service; the batch path reached completion once a single managed-dependency
role grant was in place.}
\label{tab:capabilities}
\end{table}

All four are live end-to-end. The grounded retrieval is worth a sentence of its
own, because it is the one that touches citation integrity: it queries a Vertex AI Search
datastore built over the engine's own literature warehouse, sixty-three documents imported
from the \code{literature.papers} table, and it returns real document identifiers, so a
grounded answer points to a passage that was actually retrieved rather than to an invented
reference. This is a different mechanism from the web-search grounding of the research
route in Section~\ref{sec:research-route}; here the corpus is our own. The fourth
capability, batch prediction, stages its input, submits a real asynchronous job, and writes
the resulting embeddings back to storage on completion. Reaching completion required one
managed dependency to be satisfied: the embedding batch is routed through a managed
delegation account that needed a single role grant to run on the project's behalf. With
that grant in place the path runs end-to-end, verified by a submitted job that completed
and returned its embeddings.

Two properties close the layer. A secret never leaves the boundary: a scan gate refuses to
commit or deploy any file carrying a secret-shaped value, reporting only a redacted
fingerprint, and it matches values rather than environment-variable names so the engine's
own key-reading code does not trip it; the route logs the capability name and the coded
outcome but never the parameters and never a token, and the cloud credential is minted at
the infrastructure boundary and used only as a bearer header. And the layer is reversible:
clearing a capability's flag returns it to a disabled state, and the retrieval datastore
can be deleted to return its cost to zero. The failable evaluation
\file{test/upgrade.test.mjs} exercises every gate, default-off, the typed-parameter check,
the ceiling trip with no token minted, the prerequisite hole with no spend, and the
secret-scan blocking then clearing; six governance dossiers in \file{upgrade/} carry the
inventory, the cost model, the threat model, an activation ledger, a reverification report,
and the gate.

\subsection{The SRI live feed as a standing operation}
\label{sec:sri-feed}

The systemic-risk feed is the engine's one continuously-updating product, and it is
built to run unattended. A daily Cloud Scheduler tick invokes
\route{POST /api/sri/cron-tick}; the kernel pulls the latest market closes from a
public source, appends them to the stored returns panel, recomputes the index in a
net-isolated step, and writes the new value to the \code{systemic\_risk.daily}
store. The recomputation reuses the same KSG transfer-entropy estimator as
\method{ksg\_te} (no surrogates), over the most-recent rolling window of the panel,
across every active directed pair; the index is the mean directed transfer entropy,
a system-wide nonlinear information-flow proxy where higher means more
interconnected. It is deliberately a connectivity index and is not the MCPFM
validation construct; the two are not conflated.

The operation is idempotent and self-healing. The trusted Node layer fetches the
data and reads or writes the data store, while the R step only ever sees an injected
panel and never the network, preserving net-isolation. A missing day is skipped
honestly rather than fabricated, and the panel's composition heals the gap on the
next run. The tick also exposes a dry mode that exercises the full read-and-compute
path with no writes, which serves as a health check. The feed handles markets that
drop out of the panel by excluding only their own directed pairs at the affected
date, so a date is a consistent constituent set rather than a stale mix: the first
value, on 2026-03-18, was an SRI of $0.00849$ over 306 directed pairs, and the
2026-06-05 value was $0.00709$ over 272 pairs, the smaller pair count reflecting
markets without data in that window. One operational rule governs the schedule: the
tick is never fired manually during the US trading session, because the fixed
schedule is itself the guard against computing on a partial bar.

\subsection{The nightly self-check loop}
\label{sec:nightly}

The kernel and feed are built to run unsupervised: the daily systemic-risk tick
already runs on a managed scheduler, and a nightly loop, installed in cron,
re-establishes the system's invariants and fails loudly when one breaks. The loop re-arms
the tower if its port is closed without ever killing a running job, runs the full
evaluation suite in place to refresh \file{evals.json}, refreshes the measured-state
and claims ledgers (a regression exits non-zero; a red result becomes a contested
pair and nothing is deleted), and finally, at low priority and doing zero compute,
renders the academic run report from that night's \file{evals.json} verbatim and
publishes it to the gallery. The report and gallery steps recompute nothing, so the
reporting layer can never re-trigger heavy work. The discipline of the loop -- real
compute under pre-registered failable bands, never hand-edited artefacts, honest
negatives preserved -- is the operational embodiment of the verification posture
that the rest of the manuscript develops; its own wiring is itself a failable
evaluation row.
